\section{Conclusion}

This paper develops a semiparametric framework in which a structural WGAN-GP
learns a latent valuation distribution while a cross-fitted orthogonal score
targets a regularized reserve. Under fixed-bandwidth overlap, single-crossing
and nuisance-rate conditions, the empirical-local and fixed-pseudo-count
calibrated estimators are asymptotically linear and first-order equivalent to
the same-kernel full-sample plug-in. These results do not establish an
efficiency gain from adversarial training. The score is locally orthogonal,
not a two-sided doubly robust score, and standalone WGAN inference requires
separate local-accuracy conditions.

The confidence-set construction preserves the nonlinear score and its
candidate-dependent variance. It is distinct from the Wald interval and from
the conservative interval envelope of its accepted values. The full set may
be disconnected; retaining all components preserves the test-inversion
definition. An envelope fills rejected gaps and cannot reduce coverage, but
it can increase length and does not repair a lack of information. Reporting
the original components, disconnection and boundary flags, and the envelope
separately makes this distinction observable.

The independent validation provides limited rather than conclusive support
for a small-sample coverage improvement. At $n=500$, calibrated-GAN and
empirical-local coverage rises from 91 to 93 percent under score inversion;
at $n=2000$, the constructions perform similarly near the nominal level.
Two small-sample replications produce low-end accepted components associated
with highly concentrated score variance. A sensitivity check that fixes the
variance at the estimated root yields connected sets in the replay, but
reduces small-sample coverage. The adopted response is therefore a minimal
reporting safeguard, not another tuned estimator: retain the complete set,
flag its shape and search-boundary contact, and provide an explicitly
conservative envelope when a single interval is needed.

The archived experiments explain why local accuracy and information matter:
good global distributional fit can coexist with reserve error, weakly observed
tails, and inaccurate intervals. They motivate empirical local moments and
calibration but do not demonstrate a stable reserve-estimation gain from
the GAN. Their shrinking-bandwidth and residual-bias procedures remain
diagnostic extensions, not substitutes for a theorem about the exact reserve.
The strongest current formal claim is fixed-bandwidth inference under the
stated assumptions. Further work must assess coverage across additional
data-generating processes, justify any exact-reserve bandwidth rule, and
establish a suitable empirical application without confusing price, bid and
valuation data.
